\documentclass[a4paper,10pt]{article}
\usepackage[utf8]{inputenc}
\usepackage[margin=1.5cm, bottom=2cm]{geometry}
\usepackage{fancyhdr}
\usepackage[dvipsnames]{xcolor}
\usepackage{enumitem}
\usepackage{xcolor}
\usepackage[table]{xcolor}
\usepackage{makeidx}
\usepackage[most]{tcolorbox}
\newtcolorbox{osservazioni}{
  colback=gray!8,    % sfondo grigio chiaro
  colframe=white,      % nessun bordo
  arc=2mm,            % angoli smussati
  boxrule=0pt,        % spessore bordo 0
  fontupper=\footnotesize,
  left=2mm, right=2mm, top=2mm, bottom=2mm,
  before skip=7pt,   % niente spazio prima
  after skip=15pt     % niente spazio dopo
}
\newtcolorbox{osservazioni2}{
  colback=gray!8,    % sfondo grigio chiaro
  colframe=white,      % nessun bordo
  arc=2mm,            % angoli smussati
  boxrule=0pt,        % spessore bordo 0
  fontupper=\small,
  left=5mm, right=5mm, top=3mm, bottom=3mm,
  before skip=0pt,   % niente spazio prima
  after skip=30pt     % niente spazio dopo
}
\pagestyle{fancy}
\usepackage{tikz}
\usetikzlibrary{automata, positioning}
\pagestyle{fancy}

\usepackage{makeidx}

%%%%%%%%%%%%%%%%%%%
\newcommand{\EsercitazioneNum}{11}
\newcommand{\Data}{28/11/2025}
\newcommand{\DocType}{Esercitazione} % o "Eserciziario"
\newcommand{\FullTitle}{\DocType\ \EsercitazioneNum}

\newcommand{\Header}{
\noindent
  \small \textbf{Computabilità, Complessità e Logica - a.a. 2025/26} \hfill \textbf{\FullTitle} \\[0.2cm]
  \scriptsize Matteo Liotta, \texttt{matteo.liotta@studenti.units.it} \hfill \Data \\[0.2cm]
  \noindent\makebox[\linewidth]{\rule{\linewidth}{0.4pt}} \\[0.3cm]}

\newcommand{\NewPageHeader}{
  \vfill
  \noindent\makebox[\linewidth]{\rule{\linewidth}{0.4pt}}
  \newpage
  {\Header}
}
%%%%%%%%%%%%%%%%%%%

% Numero di pagina in basso al centro
\fancyhf{}
\fancyfoot[C] \thepage
\renewcommand{\footrulewidth}{0pt} % niente linea in basso
\renewcommand{\headrulewidth}{0pt} % niente linea in alto

\begin{document}

\footnotesize  % font generale più piccolo
\Header\\[0.7cm]


%%%%%%%%%%%%%%%% FOGLIO 2

    
    %% INDEX
    \addcontentsline{toc}{section}{Esercitazione 11: \textit{CNF e Algoritmi di Soddisfacibilità}}
    
    \noindent\Large \textbf{Esercitazione 11}: \textit{CNF e Algoritmi di Soddisfacibilità (DP, DPLL)} 
    \begin{osservazioni}
        \textbf{\textit{Nota}}. Sono indicati in \textcolor{ForestGreen}{\textit{verde}} gli esercizi a sfondo teorico, proposti per rafforzare le idee alla base dei metodi utilizzati. Sono invece indicati in \textcolor{blue}{\textit{blu}} gli esercizi che non saranno trattati durante l'esercitazione, bensì lasciati come strumento aggiuntivo di esercitazione individuale. Resta comunque assoluta disponibilità per la loro risoluzione negli incontri successivi, qualora necessario. 
    \end{osservazioni}
    
    {\normalsize
    \begin{itemize}[label=, leftmargin=0pt]
    
    %\item \textbf{\underline{Formule booleane cont'd.}}

    %% INDEX
    %\addcontentsline{toc}{subsection}{Esercizio 4.}
    %%
    %\item \textbf{\textcolor{black}{Esercizio 4.} (**)}\\
    %Si considerino le seguenti formule booleane:    
    %    $$\Big(((A \rightarrow (B \land \lnot C)) \lor (\overline A \lor B)) \rightarrow (\lnot A \lor C)\Big) \lor (B \land \lnot A) $$
    %    \noindent e
    %    $$\phi_2 = \big((\lnot A \lor B) \land (\lnot C \lor B)\big) \rightarrow ((C \lor B) \land A)$$
    
    
    %Si indichi, giustificando il motivo, se la prima è equivalente alla seconda.\\[0.2cm]

    
    \item \textbf{\underline{Formule booleane in forma CNF}}
    
    %% INDEX
    \addcontentsline{toc}{subsection}{Esercizio 7.} 
    %%
    \item \textbf{\textcolor{ForestGreen}{Esercizio 7.} (**)}\\ 	
    Si consideri la seguente formula booleana:
    $$\phi:= l_1\lor l_2 \lor l_3 \lor l_4$$
    Considerata allora la scrittura:
    $$\phi':=\phi_1\land \phi_2 = (X\lor l_1\lor l_2)\land(\overline{X} \lor l_3\lor l_4)$$
    Data una assegnazione $\alpha$ tale che $\alpha \models \phi$, si indichino tutte le assegnazioni che modellano $\phi'$.


    %% INDEX
    \addcontentsline{toc}{subsection}{Esercizio 8.} 
    %%
    \item \textbf{\textcolor{blue}{Esercizio 8.} (**)}\\ 	
    Si consideri la seguente formula booleana:

    $$\phi:= \bigg(P \rightarrow (Q\lor R)\bigg)\rightarrow \bigg( (P\rightarrow Q) \lor (P\rightarrow Q) \bigg)$$
    
    Si scriva $\phi$ in forma congiuntiva di clausole CNF.
    
    %% INDEX
    \addcontentsline{toc}{subsection}{Esercizio 9.} 
    %%
    \item \textbf{\textcolor{black}{Esercizio 9.} (**)}\\ 	
    Si consideri la seguente formula booleana:

    $$\phi:=(\ \lnot(X_1\lor X_2)\ \lor \ (X_3 \lor X_4)\  )\ \land\ X_5$$
    
    Si scriva $\phi$ in forma congiuntiva di clausole CNF.\\[0.2cm]

    \textbf{\underline{Algoritmi di Soddisfacibilità: \textit{Davis-Putnam} (DP)}}
    
    %% INDEX
    \addcontentsline{toc}{subsection}{Esercizio 10.} 
    %%
    \item \textbf{Esercizio 10 (*)}\\ 	
    Si consideri il seguente insieme di clausole:
    $$\{Q,\lnot R\}, \{R,\lnot S\},\{S,\lnot T\},\{T,\lnot U\},\{U,\lnot W\},\{W,\lnot Q\}$$
    
    \begin{enumerate}
        \item Si indichi se la formula booleana $\phi$ corrispondente è soddisfacibile utilizzando l'algoritmo di Davis-Putnam (DP);
        \item Se soddisfacibile, si proponga una valutazione delle  $\alpha$ tale per cui $\alpha\models\phi$.
    \end{enumerate}

    %% INDEX
    \addcontentsline{toc}{subsection}{Esercizio 11.} 
    %%
    \item \textbf{\textcolor{black}{Esercizio 11.} (**)}\\ 	
    Si consideri il seguente insieme di clausole:
    $$\{\lnot A,B\}, \{\lnot B,C\}, \{\lnot C,D\},\{\lnot D,A\},\{\lnot A, \lnot C\},\{B, C\}$$
    \begin{enumerate}
        \item Si indichi se la formula booleana $\phi$ corrispondente è soddisfacibile utilizzando l'algoritmo di Davis-Putnam (DP);
        \item Se soddisfacibile, si proponga una valutazione delle  $\alpha$ tale per cui $\alpha\models\phi$.\\[0.2cm]
    \end{enumerate}

    \NewPageHeader
    
    %% INDEX
    \addcontentsline{toc}{subsection}{Esercizio 12.} 
    %%
    \item \textbf{\textcolor{blue}{Esercizio 12.} (**)}\\ 	
    Si consideri il seguente insieme di clausole:
    $$\{A, B, C\}, \{\lnot A, \lnot B\}, \{\lnot B, \lnot C\}, \{A, \lnot C\}, \{\lnot A, D\}, \{B, \lnot D\}$$
    \begin{enumerate}
        \item Si indichi se la formula booleana $\phi$ corrispondente è soddisfacibile utilizzando l'algoritmo di Davis-Putnam (DP);
        \item Se soddisfacibile, si proponga una valutazione delle  $\alpha$ tale per cui $\alpha\models\phi$.\\[0.2cm]
    \end{enumerate}

    %% INDEX
    \addcontentsline{toc}{subsection}{Esercizio 13.} 
    %%
    \item \textbf{\textcolor{blue}{Esercizio 13.} (**)}\\ 	
    Si consideri il seguente insieme di clausole:
    $$\{P, Q, S\}, \{P, Q, \lnot S\}, \{P, \lnot Q\}, \{\lnot P, R\}, \{\lnot P, \lnot R\}$$
    \begin{enumerate}
        \item Si indichi se la formula booleana $\phi$ corrispondente è soddisfacibile utilizzando l'algoritmo di Davis-Putnam (DP);
        \item Se soddisfacibile, si proponga una valutazione delle  $\alpha$ tale per cui $\alpha\models\phi$.\\[0.2cm]
    \end{enumerate}

    %% INDEX
    \addcontentsline{toc}{subsection}{Esercizio 14.} 
    %%
    \item \textbf{\textcolor{blue}{Esercizio 14.} (**)}\\ 	
    Si consideri il seguente insieme di clausole:
    $$\{A, B\}, \{\lnot A, C\}, \{\lnot B, D\}, \{\lnot C, \lnot D\}, \{C, \lnot B\}, \{A, \lnot D\}$$
    \begin{enumerate}
        \item Si indichi se la formula booleana $\phi$ corrispondente è soddisfacibile utilizzando l'algoritmo di Davis-Putnam (DP);
        \item Se soddisfacibile, si proponga una valutazione delle  $\alpha$ tale per cui $\alpha\models\phi$.\\[0.2cm]
    \end{enumerate}

    \textbf{\underline{Algoritmi di Soddisfacibilità: \textit{Davis-Putnam-Logemann-Loveland} (DPLL)}}

    %% INDEX
    \addcontentsline{toc}{subsection}{Esercizio 15.} 
    %%
    \item \textbf{\textcolor{ForestGreen}{Esercizio 15.} (**)}\\ 	
    Si consideri la formula booleana in forma CNF:
    $$\phi = (P)\land(\overline P \lor R)\land (\overline R \lor Q)$$
    \begin{enumerate}
        \item Si indichi se la formula booleana $\phi$ corrispondente è soddisfacibile utilizzando l'algoritmo di Davis-Putnam-Logemann-Loveland (DPLL);
        \item Se soddisfacibile, si proponga una valutazione delle  $\alpha$ tale per cui $\alpha\models\phi$.
    \end{enumerate}

    %% INDEX
    \addcontentsline{toc}{subsection}{Esercizio 16.} 
    %%
    \item \textbf{Esercizio 16. (**)}\\ 	
    Si consideri il seguente insieme di clausole:
    $$\{A, B, \lnot C\}, \{\lnot A, C, D\}, \{\lnot B, C, \lnot D\}, \{\lnot C, D, E\}, \{\lnot D, \lnot E, B\}, \{\lnot A, \lnot B, E\}$$
    \begin{enumerate}
        \item Si indichi se la formula booleana $\phi$ corrispondente è soddisfacibile utilizzando l'algoritmo di Davis-Putnam-Logemann-Loveland (DPLL);
        \item Se soddisfacibile, si proponga una valutazione delle  $\alpha$ tale per cui $\alpha\models\phi$.
    \end{enumerate}

    %% INDEX
    \addcontentsline{toc}{subsection}{Esercizio 17.} 
    %%
    \item \textbf{Esercizio 17. (**)}\\ 	
    Si consideri il seguente insieme di clausole:
    $$\{A, B\}, \{A, \neg B\}, \{\neg A, C\}, \{\neg A, \neg C\}$$
    \begin{enumerate}
        \item Si indichi se la formula booleana $\phi$ corrispondente è soddisfacibile utilizzando l'algoritmo di Davis-Putnam-Logemann-Loveland (DPLL);
        \item Se soddisfacibile, si proponga una valutazione delle  $\alpha$ tale per cui $\alpha\models\phi$.\\[0.2cm]
    \end{enumerate}

    \NewPageHeader

    %% INDEX
    \addcontentsline{toc}{subsection}{Esercizio 18.} 
    %%
    \item \textbf{\textcolor{blue}{Esercizio 18.} (**)}\\ 	
    Si consideri il seguente insieme di clausole:
    $$\{P, Q, R\}, \{\lnot P, Q, S\}, \{\lnot Q, R, S\}, \{\lnot R, \lnot S, Q\}, \{\lnot P, \lnot Q, \lnot R\}, \{P, \lnot Q, \lnot S\}$$
    \begin{enumerate}
        \item Si indichi se la formula booleana $\phi$ corrispondente è soddisfacibile utilizzando l'algoritmo di Davis-Putnam-Logemann-Loveland (DPLL);
        \item Se soddisfacibile, si proponga una valutazione delle  $\alpha$ tale per cui $\alpha\models\phi$.
    \end{enumerate}

    %% INDEX
    \addcontentsline{toc}{subsection}{Esercizio 19.} 
    %%
    \item \textbf{\textcolor{blue}{Esercizio 19.} (**)}\\ 	
    Si consideri il seguente insieme di clausole:
    $$\{X, Y, \lnot Z\}, \{\lnot X, Z, U\}, \{\lnot Y, Z, V\}, \{\lnot Z, \lnot U, W\}, \{\lnot V, \lnot W, Y\}, \{\lnot X, \lnot V, W\}, \{X, \lnot Y, V\}$$
    \begin{enumerate}
        \item Si indichi se la formula booleana $\phi$ corrispondente è soddisfacibile utilizzando l'algoritmo di Davis-Putnam-Logemann-Loveland (DPLL);
        \item Se soddisfacibile, si proponga una valutazione delle  $\alpha$ tale per cui $\alpha\models\phi$.
    \end{enumerate}

    %% INDEX
    \addcontentsline{toc}{subsection}{Esercizio 20.} 
    %%
    \item \textbf{\textcolor{blue}{Esercizio 20.} (**)}\\ 	
    Si consideri il seguente insieme di clausole:
    $$\{A, B, C, D, \overline E\}, \{C, \overline B\}, \{D, \overline A\}, \{B, \overline D\}, \{C, A, \overline E\}, \{A, B, C, D, E\},$$$$\{C, \overline D\}, \{\overline D, \overline A, B\}, \{E, \overline C, D\}, \{C, A, \overline E\}$$
    \begin{enumerate}
        \item Si indichi se la formula booleana $\phi$ corrispondente è soddisfacibile utilizzando \textbf{sia} l'algoritmo di Davis-Putnam-Logemann-Loveland (DPLL) che l'algoritmo di Davis-Putnam (DP);
        \item Se soddisfacibile, si proponga una valutazione delle  $\alpha$ tale per cui $\alpha\models\phi$.
    \end{enumerate}

    
\end{itemize}
\vfill
\noindent\makebox[\linewidth]{\rule{\linewidth}{0.4pt}}
\end{document}